\newpage

% ######################################################################
% PART IX: EXTENDED Q&A, TROUBLESHOOTING, TUTORIAL
% ######################################################################
\part{Extended Q\&A, Troubleshooting \& Tutorial}

% ######################################################################
% SECTION 40: 30 MORE INTERVIEW QUESTIONS
% ######################################################################
\section{30 More Interview Questions --- Advanced \& Career}

\subsection{Advanced LLM Topics}

\begin{interviewbox}
\textbf{Q49: ``Fine-tuning vs RAG --- kab kya?''}

\textbf{Answer:} RAG: \textbf{knowledge injection} --- facts/docs jo badalte rehte hain ya private hain. Fine-tuning: \textbf{behavior/style} --- format, tone, domain-specific response pattern. Rule of thumb: RAG pehle try karo (cheap, updateable, no training), fine-tuning tab jab style/behavior problem ho jo prompting se na solve ho. Is project mein RAG use kiya (knowledge) --- fine-tuning zaroori nahi thi.
\end{interviewbox}

\begin{interviewbox}
\textbf{Q50: ``RLHF kya hai?''}

\textbf{Answer:} Reinforcement Learning from Human Feedback: (1) Model ko responses generate karwao; (2) Humans rank karein; (3) Reward model train karo; (4) Model ko RL se align karo. Ye LLM ko \textbf{helpful, harmless, honest} banata hai. Pre-trained model (next-word) vs aligned model (useful assistant) --- API models aligned hote hain.
\end{interviewbox}

\begin{interviewbox}
\textbf{Q51: ``Quantization kya hai?''}

\textbf{Answer:} Model weights ki precision kam karna: FP32 (4 bytes/weight) $\rightarrow$ FP16/INT8/INT4. 70B model FP16 = 140GB, INT4 = 35GB --- RAM requirement 4x kam. Trade-off: thodi accuracy loss, bada memory/speed gain. Is project mein models provider side hain (Groq), par RAG embeddings local chalti hain --- quantization wahan relevant hota.
\end{interviewbox}

\begin{interviewbox}
\textbf{Q52: ``Context window ko extend kaise karein?''}

\textbf{Answer:} 4 techniques: (1) \textbf{Positional encoding} improvements (RoPE, ALiBi) --- extrapolate karta hai; (2) \textbf{Sliding window} attention; (3) \textbf{Retrieval} --- pura context dena hi mat, relevant parts do (is project ki strategy); (4) \textbf{Hierarchical} summaries. ``Is project mein context compression: analyst dense brief banata hai, writer ko sirf wo + top-8 sources jate hain.''
\end{interviewbox}

\begin{interviewbox}
\textbf{Q53: ``Embedding model vs LLM mein difference?''}

\textbf{Answer:} Embedding model: text $\rightarrow$ vector (retrieval ke liye, encoder-only, fast). LLM: text $\rightarrow$ text (generation, decoder-only, heavy). Is project mein dono: all-MiniLM-L6-v2 embeddings (ChromaDB ke liye) + GPT-OSS-120B generation.
\end{interviewbox}

\begin{interviewbox}
\textbf{Q54: ``Streaming responses kyun important?''}

\textbf{Answer:} User perception: 176s ka request bina feedback = user sochta hai system hang hai. Streaming token-by-token output dikhata hai --- perceived latency kam. Is project mein Streamlit agent-status updates real-time dikhata hai (graph.stream events), aur MCP/API invoke() mein final result hi milta hai.
\end{interviewbox}

\subsection{MLOps \& Production}

\begin{interviewbox}
\textbf{Q55: ``Prompt versioning kya hai? Kya version control karte ho?''}

\textbf{Answer:} Prompts bhi code hain --- git mein version karo, change log rakho, aur A/B test karo (same prompt ke 2 versions benchmark par). Is project mein prompts constants hain (agents/*.py) --- production mein separate prompt registry + config-driven banata.
\end{interviewbox}

\begin{interviewbox}
\textbf{Q56: ``Observability: kya track karte ho?''}

\textbf{Answer:} (1) \textbf{Per-node latency} --- researcher kitna, writer kitna; (2) \textbf{Token/cost} per run; (3) \textbf{Retry/failure rates} --- API errors; (4) \textbf{Quality metrics} --- benchmark scores over time; (5) \textbf{Tracing} --- LangSmith/Langfuse se pura run graph. Is project mein Django DB runs store karta hai + benchmark history.
\end{interviewbox}

\begin{interviewbox}
\textbf{Q57: ``CI/CD is project mein kaise lagate?''}

\textbf{Answer:} GitHub Actions pipeline: (1) \textbf{Lint/format} (ruff/black); (2) \textbf{pytest} --- 16 tests; (3) \textbf{Build} --- Docker images; (4) \textbf{Deploy} --- Streamlit Cloud/Django server. Quality gate: tests fail = deploy nahi. Ye standard practice hai jo project par seedha apply hoti hai.
\end{interviewbox}

\begin{interviewbox}
\textbf{Q58: ``Feature flag kya hota hai?''}

\textbf{Answer:} Code deploy without enabling feature --- runtime config se on/off. Use: gradual rollout, A/B testing, instant rollback. Is project mein example: naya analyst prompt flag ke peeche --- 10\% traffic par enable karke benchmark compare karo, phir full rollout.
\end{interviewbox}

\begin{interviewbox}
\textbf{Q59: ``Cost optimization kaise karte?''}

\textbf{Answer:} (1) \textbf{Caching} --- same topic ka result cache; (2) \textbf{Model tiering} --- simple tasks chhota model, complex bada; (3) \textbf{Context discipline} --- kam tokens; (4) \textbf{Batching} --- multiple requests ek call mein; (5) \textbf{Local models} --- embeddings free. Is project mein: embeddings local free, Groq cheap, sirf Tavily paid --- cache search results.
\end{interviewbox}

\begin{interviewbox}
\textbf{Q60: ``A/B testing LLM features kaise?''}

\textbf{Answer:} (1) Traffic split (50/50); (2) Har version ke outputs \textbf{LLM judge} se score; (3) Metrics: depth, verifiability, user feedback; (4) \textbf{Statistical significance} --- p-value, sample size (ek run ka verdict meaningless, jaise is project ke benchmark mein dikha). Ye hi is project ke benchmark infra ka production use case hai.
\end{interviewbox}

\subsection{Career \& Team}

\begin{interviewbox}
\textbf{Q61: ``Apni strength/weakness batao?''}

\textbf{Answer:} \textbf{Strength:} Problem solving with measurement --- main benchmark/data se decisions leta hoon (is project ke examples se prove karo). \textbf{Weakness:} Kabhi kabhi detail mein itna chala jata hoon ki high-level deadline miss hoti --- maine time-boxing seekhi (sprint planning, priorities). Weakness ko \textbf{improvement plan} ke saath bolna zaroori hai.
\end{interviewbox}

\begin{interviewbox}
\textbf{Q62: ``5 saal baad khud ko kahan dekhte ho?''}

\textbf{Answer:} (Honest + growth-oriented) ``AI engineering mein deep expertise --- complex production systems (agents, RAG, evaluation infra) design karna. Pehle 2 saal: strong foundation (system design, ML fundamentals, production engineering). Aage: team lead / architect role jahan main systems design karun aur juniors ko mentor karun. Ye project usi direction ka pehla step hai.''
\end{interviewbox}

\begin{interviewbox}
\textbf{Q63: ``Salary expectation?''}

\textbf{Answer:} Research karke bolo (Glassdoor/levels.fyi, role+location basis). Range do: ``Maine is role ke liye 12-15 LPA ka range dekha hai, experience basis par. Main flexible hoon agar growth + learning acha hai.'' Kisi number se pehle interviewer ko bolne do (``aapke budget ke hisaab se kya range hai?'') --- negotiation basics.
\end{interviewbox}

\begin{interviewbox}
\textbf{Q64: ``Do you have any questions for us?''}

\textbf{Answer:} 3 acha sawaal: (1) ``Team mein AI features kaise ship hote hain --- evaluation/benchmarking process kya hai?'' (2) ``Is role mein pehle 90 din mein kya achieve karna expected hai?'' (3) ``Team kaise kaam karti hai --- code review, on-call, sprint structure?'' Ye questions aapki \textbf{engineering maturity} dikhate hain.
\end{interviewbox}

\begin{interviewbox}
\textbf{Q65: ``Why NIT Hamirpur / dual degree CSE?''}

\textbf{Answer:} (Personal + honest) ``NIT Hamirpur ka CSE program strong theoretical foundation deta hai --- DSA, OS, DBMS, ML core courses. Dual degree (B.Tech + M.Tech) ne research exposure diya --- thesis/project work mein deep-dive ka habit bana. Is project mein wo hi habit dikhti hai --- end-to-end system, benchmark se validate.''
\end{interviewbox}

\begin{interviewbox}
\textbf{Q66: ``College projects / internships batao?''}

\textbf{Answer:} (STAR format, 2-3 projects ready rakho) Is project ko main story banao + 1-2 aur (web dev, ML, hackathon). Har project mein: problem $\rightarrow$ approach $\rightarrow$ tech stack $\rightarrow$ result (numbers ke saath) $\rightarrow$ lessons.
\end{interviewbox}

\begin{interviewbox}
\textbf{Q67: ``Coding round ke liye DSA preparation kaise?''}

\textbf{Answer:} (Interview context) ``Arrays/strings/hashmaps $\rightarrow$ two pointers, sliding window, binary search $\rightarrow$ trees/graphs $\rightarrow$ DP. 300+ problems, pattern-based (not count-based). Weekly contests + mock interviews. Ye project DSA se alag skill hai --- dono parallel prepare karna zaroori.''
\end{interviewbox}

\begin{interviewbox}
\textbf{Q68: ``Aapne ye project kaise time management kiya?''}

\textbf{Answer:} (Honest) ``Modules alag the (backend, agents, UI, benchmark) --- har weekend ek module. Pehle working end-to-end (thoda raw), phir polish. Testing last mein nahi, har module ke saath. Failures (benchmark loss) ko feature ke roop mein treat kiya --- measurement infra banaya. Realistic: 4-6 weeks, evenings + weekends.''
\end{interviewbox}

% ######################################################################
% SECTION 41: TROUBLESHOOTING GUIDE
% ######################################################################
\section{Troubleshooting Guide --- Common Problems \& Fixes}

\begin{longtable}{p{5.2cm} p{4.2cm} p{4.6cm}}
\toprule
\textbf{Symptom} & \textbf{Cause} & \textbf{Fix} \\
\midrule
No module named django & venv activate fail / deps missing & .venv/Scripts/python.exe + pip install -r \\
FastMCP unexpected keyword & fastmcp version mismatch & instructions= instead of description= \\
Model 404 / bad response & MODEL\_NAME unavailable & GET /v1/models, live test, update config \\
Port 8000 already in use & Another service & runserver 8001 + BACKEND\_API\_URL \\
Emoji print error (Windows) & Console encoding & PYTHONIOENCODING=utf-8 \\
Test time frozen / ties & Django test env & order by -created\_at, -id \\
LaTeX Missing \$ & Underscore in caption/text & Escape \_ \\
Streamlit secrets missing & Cloud secrets not set & Add keys in Streamlit settings \\
ChromaDB empty results & Collection empty / model mismatch & Check stats endpoint, same embedding model \\
Slow first run & Model download / warmup & Preload embeddings, cache \\
\bottomrule
\end{longtable}

% ######################################################################
% SECTION 42: 30-MINUTE TUTORIAL
% ######################################################################
\section{30-Minute Tutorial --- Clone Se Running Tak}

\begin{codestorybox}
\textbf{Minute 0-5: Setup}
\begin{lstlisting}[style=bashstyle]
git clone <repo-url>
cd <project>
python -m venv .venv
# Windows: .venv\Scripts\activate   |   Linux/Mac: source .venv/bin/activate
pip install -r requirements.txt
cp .env.example .env   # keys daalo
\end{lstlisting}

\textbf{Minute 5-10: Backend}
\begin{lstlisting}[style=bashstyle]
cd backend
python manage.py migrate
python manage.py runserver 8001
# verify: curl http://127.0.0.1:8001/api/health/
\end{lstlisting}

\textbf{Minute 10-15: UI}
\begin{lstlisting}[style=bashstyle]
# naya terminal
streamlit run app.py
# browser: http://localhost:8501
\end{lstlisting}

\textbf{Minute 15-20: First Research}
\begin{itemize}
\item Topic type karo (jaise ``AI in healthcare 2026'')
\item Run dabao --- agent statuses live dikhenge
\item Report + download + PDF upload try karo
\end{itemize}

\textbf{Minute 20-25: Benchmark Lab}
\begin{itemize}
\item Sidebar mein ``Benchmark'' page
\item Topic do, Run --- single vs multi-agent comparison dekho
\item History charts dekho
\end{itemize}

\textbf{Minute 25-30: MCP + Tests}
\begin{lstlisting}[style=bashstyle]
pytest                 # 16 tests, ~9s
python mcp_server.py   # MCP server
\end{lstlisting}

\textbf{Done!} Interview demo ke liye ye script yaad rakho --- 30 min mein poora system dikhaya.
\end{codestorybox}

% ######################################################################
% SECTION 43: PROJECT FILE MAP
% ######################################################################
\section{Project File Map --- Har File Ka Kaam}

\begin{longtable}{p{5.2cm} p{9.2cm}}
\toprule
\textbf{File} & \textbf{Kaam} \\
\midrule
app.py & Streamlit main page: UI, session state, streaming pipeline, download \\
pages/1\_Benchmark.py & Benchmark Lab: run benchmark + history charts \\
mcp\_server.py & MCP server: 5 tools + resource for AI clients \\
config/setting.py & Central config: env vars, model, chunk settings, API keys \\
state/schema.py & ResearchState TypedDict --- pipeline data contract \\
graph/workflow.py & LangGraph: nodes, edges, conditional loop, compiled singleton \\
agents/researcher.py & Search: sub-queries, Tavily, RAG retrieval \\
agents/analyst.py & Synthesis: themes, contradictions, gaps \\
agents/fact\_checker.py & Verification: closed-world claims check, VERDICT \\
agents/writer.py & Report generation + memory store \\
rag/vector\_store.py & ChromaDB: ingest, retrieve, stats, clear \\
benchmark/evaluator.py & Baseline + evaluation (LLM judge + heuristics) \\
backend/manage.py & Django entry point \\
backend/orchestrator\_backend/ & Django settings, urls, wsgi/asgi \\
backend/api/models.py & ResearchSession + BenchmarkEvaluation \\
backend/api/serializers.py & DRF serializers + validation \\
backend/api/views.py & 6 API endpoints \\
backend/api/urls.py & URL routing \\
backend/api/tests/ & 16 pytest tests (auth, research, benchmark) \\
backend/conftest.py & Mock fixtures for tests \\
requirements.txt & Python dependencies \\
pytest.ini & pytest + Django config \\
.env.example & Environment template \\
README.md & Project docs + API reference \\
Multi\_Agent\_Master\_Notes.tex & Ye document (source) \\
\bottomrule
\end{longtable}

\begin{memorybox}
\textbf{Final revision tip:} Interview se pehle: (1) 30-min tutorial 2 baar practice karo; (2) 10-15 Q\&A khud bol ke practice karo (voice mein); (3) Project ke 3 numbers yaad rakho: 16 tests, 4 agents, 3 interfaces; (4) Ye document khol ke rakho --- koi bhi question, kahin bhi jump karo. \textbf{All the best!}
\end{memorybox}
